% SPDX-License-Identifier: GPL-3.0-or-later
% Copyright (C) 2026 Alessandro Greco
\section*{Глава \ref{cap:11-stream}: потоки, SENDME и управление перегрузкой}

\begin{soluzione}{es:11-begin}
Имя хоста должно быть в нижнем регистре: \texttt{example.com:443}, затем нулевой байт, то есть \texttt{6578616d706c652e636f6d3a343433 00}, после чего FLAGS с установленным битом 1 для разрешения IPv6: \texttt{00000001}.
\end{soluzione}

\begin{soluzione}{es:11-connected}
IPv4-адрес \texttt{c0000201} — это 192.0.2.1; TTL \texttt{0000012c} равен 300 секундам. Клиент может игнорировать адрес. Если он сохраняет его, нельзя повторно использовать значение между цепями: выходной ретранслятор мог передать ложные сведения или вернуть уникальный адрес, чтобы впоследствии распознать клиента.
\end{soluzione}

\begin{soluzione}{es:11-optimistic}
Без опережающей отправки данных (optimistic data) клиент ждёт \texttt{CONNECTED} перед отправкой запроса, что добавляет ещё один обмен туда и обратно примерно в 600\,мс к операции, и без того требующей как минимум двух RTT. Если соединение завершается неудачей, данные в очереди нужно отбросить, поскольку назначения, способного их получить, нет. Их нельзя передавать через другое соединение.
\end{soluzione}

\begin{soluzione}{es:11-velocita}
Поток может иметь в полёте не более 500 ячеек, то есть примерно $500 \times 498 \approx 249$\,KB на RTT. При RTT 400\,мс это около 620\,KB/s. Даже при множестве потоков цепь ограничена 1\,000 ячейками на RTT, то есть примерно 1,2\,MB/s. При управлении перегрузкой окно начинается всего со 124 ячеек, около 62\,KB, но растёт до тех пор, пока измерения RTT не покажут перегрузку, без фиксированного верхнего предела.
\end{soluzione}

\begin{soluzione}{es:11-sendme-falso}
Вредоносный клиент открывает много потоков к серверу, отправляющему большой объём данных, возможно к собственному серверу атакующего, почти ничего не читает и многократно отправляет поддельные сообщения \texttt{SENDME}. Выходной ретранслятор продолжает передачу, а ячейки накапливаются в очередях входного и промежуточного ретрансляторов, а также клиента, расходуя память ретрансляторов. Аутентифицированное сообщение \texttt{SENDME} требует доказательства, что клиент получил ячейку, вызвавшую это сообщение. Если бы содержимое ячейки было предсказуемым, например при загрузке знакомого файла, клиент мог бы вычислить дайджест, не получая саму ячейку; случайные байты заполнения предотвращают такой обход.
\end{soluzione}

\begin{soluzione}{es:11-resolve}
Общий план: отправить \texttt{RESOLVE} с новым StreamID и именем хоста, завершённым нулём, затем интерпретировать \texttt{RESOLVED} как последовательность кортежей, содержащих тип, длину, значение и TTL. Для несуществующего имени хоста выходной ретранслятор отправляет запись с кодом ошибки (\texttt{0xF0} или \texttt{0xF1}), содержимое которой следует игнорировать.
\end{soluzione}

\begin{soluzione}{es:11-policy}
Выходной ретранслятор отвечает \texttt{END} с причиной 4 (\texttt{EXITPOLICY}), после которой идут адрес назначения, четыре байта для IPv4 и четырёхбайтовый TTL. Клиент может использовать эту информацию, чтобы не повторять попытку через тот же выходной ретранслятор для данного назначения.
\end{soluzione}

\begin{soluzione}{es:11-rtt}
Мы измеряем время, за которое \texttt{SENDME} достигает выходного ретранслятора, плюс время, необходимое разблокированному выходному ретранслятору для отправки дополнительных ячеек и их доставки клиенту. Когда выходной ретранслятор ждёт \texttt{SENDME}, это приближённо соответствует RTT цепи с точки зрения клиента. Вторая одновременная загрузка удлиняет очереди и повышает измеренные значения — именно этот сигнал используется механизмом управления перегрузкой для уменьшения окна.
\end{soluzione}

\section*{Глава \ref{cap:12-directory}: система каталогов}

\begin{soluzione}{es:12-validita}
Интервал голосования равен одному часу: от \texttt{valid-after} до \texttt{fresh-until}. Клиент выбирает случайный момент между $11{:}00 + \tfrac34 \cdot 60' = 11{:}45$ и $11{:}45 + \tfrac78 \cdot 75' \approx 12{:}50$. Если к 13:00 нового консенсуса всё ещё нет, существующий формально истёк. Однако клиент продолжает использовать его как \emph{reasonably live}, то есть в течение 24 часов после окончания срока действия, согласно \file{REASONABLY\_LIVE\_TIME} в \file{networkstatus.c}. Он может продолжать строить цепи, многократно пытаясь получить более новый консенсус. Примерно через 24 часа после истечения клиент перестаёт полагаться на старый документ.
\end{soluzione}

\begin{soluzione}{es:12-flag}
Вероятнее всего, \texttt{Running} и \texttt{Valid}, а также \texttt{Fast}, поскольку 500\,KB/s больше 100\,KB/s. Возможно \texttt{V2Dir}, если ретранслятор также работает как кэш каталогов. Но не \texttt{Stable}, \texttt{Guard} или \texttt{HSDir}, для которых требуются дни или недели работы. И не \texttt{Exit}, поскольку политика выхода разрешает только порт 443, тогда как для этого флага требуется выходной трафик и на порту 80, и на 443.
\end{soluzione}

\begin{soluzione}{es:12-partizione}
Кэш отправляет выбранному клиенту двухчасовой, но всё ещё действительный консенсус, в котором отсутствуют недавно добавленные ретрансляторы. Поэтому клиент выбирает маршруты из немного иного множества, чем остальные, что потенциально делает его трафик узнаваемым. Атака ограничена трёхчасовой действительностью консенсуса, необходимостью действительной подписи центров каталогов — кэш не может придумать собственный консенсус — и случайным выбором кэшей клиентами. Загрузка микродескрипторов по дайджесту, а не по идентичности ретранслятора, не позволяет кэшу предоставить конкретному клиенту другую версию ключа ретранслятора, поскольку дайджест зафиксирован подписанным консенсусом.
\end{soluzione}

\begin{soluzione}{es:12-mediana}
При трёх измерениях медиана равна 5200. Одно завышенное измерение её не меняет. Если бы третий центр каталогов была единственной, измерявшим этот ретранслятор, система не имела бы трёх измерений, необходимых для медианы. Вместо этого консенсус использовал бы заявленную или указанную в голосах пропускную способность с \texttt{Unmeasured=1} и верхним ограничением, поэтому завышенное измерение не повлияло бы на результат.
\end{soluzione}

\begin{soluzione}{es:12-voti}
В документе голосования присутствуют строки, отсутствующие в консенсусе или отличающиеся от него. Среди примеров — значения пропускной способности \texttt{Measured} от центров измерения пропускной способности, строки \texttt{m} с поддерживаемыми методами консенсуса, флаги, назначенные конкретным центром каталогов, её \texttt{flag-thresholds} и обязательства для совместного случайного значения. Консенсус сохраняет только объединённые результаты.
\end{soluzione}

\begin{soluzione}{es:12-md}
Хэшируйте текст от строки \texttt{onion-key}, открывающей микродескриптор, до перевода строки перед следующим микродескриптором либо до конца документа, включая строки PEM-кодированного RSA-ключа. Дайджест SHA-256, закодированный base64 без завершающих символов \texttt{=}, должен совпасть со строкой \texttt{m} консенсуса.
\end{soluzione}

\begin{soluzione}{es:12-autorita}
При остановке одного центра каталогов две оставшиеся продолжают голосовать и создают консенсус с двумя подписями, чего достаточно для требуемого большинства $\lfloor 3/2 \rfloor + 1 = 2$. При остановке двух центров каталогов единственная оставшаяся не может получить большинство, и новый консенсус не создаётся. Тем не менее клиент продолжает использовать последний консенсус как \emph{reasonably live} примерно 24 часа после окончания срока действия. В лабораторной сети консенсус действует только 40 секунд, но срок допустимого использования устаревшего консенсуса всё равно равен 24 часам, поэтому нарушение не видно немедленно.
\end{soluzione}

\begin{soluzione}{es:12-ns}
Консенсус \texttt{ns} обычно содержит подписи SHA-1; спецификация для этого варианта разрешает только их. Его строка \texttt{r} дополнительно содержит полный дайджест дескриптора ретранслятора, а каждая запись ретранслятора также включает строку \texttt{p} со сводкой политики выхода.
\end{soluzione}

\section*{Глава \ref{cap:13-percorso}: выбор маршрута}

\begin{soluzione}{es:13-ordine}
К выходному ретранслятору предъявляются самые строгие требования: он должен разрешать запрошенное назначение. Выбор выходного ретранслятора первым устанавливает, что подходящий узел вообще существует. Если сначала выбрать входной узел, ограничения семейства или сети могут исключить единственный подходящий выходной ретранслятор и потребовать повторного выбора. Есть и статистическая причина: исключения применяются к последующим позициям, поэтому они должны сильнее влиять на позиции с более крупными множествами кандидатов.
\end{soluzione}

\begin{soluzione}{es:13-peso}
Как входной узел: $20000 \times 0/10000 = 0$ (\texttt{Wgd=0}). Как промежуточный узел: $20000 \times 0 = 0$ (\texttt{Wmd=0}). Как выходной узел: $20000 \times 10000/10000 = 20000$ (\texttt{Wed=10000}). При весах этого консенсуса ретранслятор используется исключительно как выходной узел.
\end{soluzione}

\begin{soluzione}{es:13-guard}
Без постоянных входных узлов вероятность компрометации превышает 50\%, когда $1 - (1 - c^2)^C > 0.5$, то есть $C > \ln 0.5 / \ln(1 - c^2)$: примерно \num{6932} цепи для $c = 1\%$ и 17 цепей для $c = 20\%$. С одним фиксированным входным узлом вероятность остаётся близкой к $c$, то есть 1\% и 20\% соответственно. Входные узлы недостаточны против очень мощного противника, но различие всё равно существенно.
\end{soluzione}

\begin{soluzione}{es:13-sybil}
Атакующий должен предоставить реальную пропускную способность: вероятности выбора зависят от пропускной способности, измеренной центрами каталогов, а не только от заявленных значений. Выходные ретрансляторы также тщательно контролируются, в том числе через \texttt{BadExit}. Ограничение в два ретранслятора на адрес вынуждает атакующего получать множество адресов. Объявление семейства не позволяет двум его ретрансляторам одновременно попасть в одну цепь; отсутствие объявления создаёт риск обнаружения и исключения центрами каталогов. Наконец, веса резервируют дефицитную выходную пропускную способность для настоящих выходных ретрансляторов, и именно там учитывается пропускная способность атакующего.
\end{soluzione}

\begin{soluzione}{es:13-campione}
Без ограничения выборки атакующий, контролирующий локальную сеть клиента, мог бы блокировать соединения ко всем честным входным узлам. Клиент продолжал бы выбирать новые входные узлы, пока не нашёл доступный, и в конце концов выбрал бы вредоносный. При ограниченной выборке клиент исчерпывает попытки и перестаёт добавлять входные узлы, возможно теряя связность вместо автоматического перехода на входной узел атакующего.
\end{soluzione}

\begin{soluzione}{es:13-famiglie}
Это практическое упражнение. Получите микродескрипторы из CollecTor вместе с соответствующими консенсусами. Для каждого ретранслятора постройте множество идентификаторов семейства через \texttt{family-ids} или взаимные списки \texttt{family} и отклоняйте кандидата, если он принадлежит к одному семейству с любым уже выбранным ретранслятором.
\end{soluzione}

\begin{soluzione}{es:13-operatori}
Это исследовательское упражнение. Крупные семейства обычно дают заметно большую долю вероятности выбора выходного узла, чем любой отдельный крупный ретранслятор. Поэтому фактический риск концентрации зависит от операторов, а не только от числа ретрансляторов. Именно этим объясняется внимание сообщества к разнообразию операторов и необъявленным семействам ретрансляторов (\cref{cap:27-casi-rete}).
\end{soluzione}

\section*{Глава \ref{cap:14-onion-services}: onion-сервисы v3}

\begin{soluzione}{es:14-indirizzo}
$32 + 2 + 1 = 35$ байт, то есть 280 бит. Base32 кодирует по пять бит на символ, поэтому требуется $280/5 = 56$ символов без символов заполнения. Случайная ошибка ввода имеет примерно вероятность $2^{-16}$ случайно пройти проверку контрольной суммы. Байт версии также должен остаться \texttt{0x03}, что делает незамеченную случайную ошибку ещё менее вероятной.
\end{soluzione}

\begin{soluzione}{es:14-periodo}
В 09:00 UTC 1 января 2026 года с начала эпохи Unix прошло \num{29454300} минут. После вычитания смещения 720 минут остаётся \num{29453580}, поэтому $\lfloor 29453580 / 1440 \rfloor = 20453$. Период \num{20453} начался в 12:00 UTC 31 декабря 2025 года, поскольку 09:00 всё ещё раньше полудня.
\end{soluzione}

\begin{soluzione}{es:14-hsdir}
Атакующему нужны ослеплённый ключ периода и, следовательно, знание адреса сервиса, а также SRV этого периода и позиции контролируемых ретрансляторов в кольце. Эти значения используются для генерации идентичностей ретрансляторов, располагающихся непосредственно после индексов сервиса. SRV становится известен только при его публикации центрами каталогов, примерно за половину суток до начала соответствующего периода. Атакующий не может подготовить ретрансляторы с точно подобранным положением за несколько дней, а новым ретрансляторам ещё требуется получить флаг \texttt{HSDir}, что требует минимум 96 часов работы.
\end{soluzione}

\begin{soluzione}{es:14-ruoli}
HSDir видит ослеплённый ключ, а не адрес сервиса, и получает запросы через анонимные цепи. Точка представления видит ключ аутентификации и цепь, приходящую от сервиса, и может подсчитывать запросы на представление. Точка встречи видит идентификатор cookie, две цепи и зашифрованный трафик между ними. Точка представления особенно полезна для оценки популярности, поскольку наблюдает каждую попытку соединения, хотя не знает соответствующий адрес. Однако сторона, уже знающая адрес, может загрузить дескриптор и установить эту связь.
\end{soluzione}

\begin{soluzione}{es:14-anonimo}
В цепи из одного перехода HSDir является первым ретранслятором и видит IP-адрес клиента. Он мог бы связать \emph{того, кто} ищет сервис, с \emph{тем, какой} сервис ищется, как минимум по ослеплённому ключу, а если адрес уже известен — и непосредственно с адресом. Отклонение неанонимных запросов вынуждает клиентов загружать дескрипторы через другие ретрансляторы.
\end{soluzione}

\begin{soluzione}{es:14-anello}
Это практическое упражнение. Когда HSDir становится больше, чем $2 \times 4$, сервис уже не публикует дескриптор на каждом. Запрос дескриптора у каждого HSDir через трёхузловую цепь должен находить его именно на HSDir, рассчитанных для текущего или следующего периода, поскольку сервис публикует дескрипторы для обоих периодов.
\end{soluzione}

\begin{soluzione}{es:14-ipt}
Значение base64 содержит число спецификаторов соединения, после которого идут сами спецификаторы в формате \texttt{EXTEND2} (\cref{sec:extend2}). Например, тип \texttt{00} задаёт адрес и ORPort, тип \texttt{02} — отпечаток RSA, тип \texttt{03} — идентичность Ed25519. Ретранслятор можно определить по ORPort, который равен 5100 плюс номер узла. Строка \texttt{auth-key} содержит сертификат типа 9, подписанный ключом подписи дескриптора и удостоверяющий ключ аутентификации сервиса в этой точке представления.
\end{soluzione}
